\subsection{Rúbricas como ayudas durante el entrenamiento}
Además de construir la recompensa, otra cuestión es cómo lograr que el modelo genere respuestas suficientemente buenas para aprender de ellas durante el entrenamiento. Zhou y sus colaboradores abordaron este problema mediante RuscaRL, un método que utiliza las rúbricas tanto para orientar la generación como para evaluar las respuestas \cite{zhou2026ruscarl}. Para ello, proporcionan distintos niveles de ayuda dentro de un mismo grupo de generaciones y reducen progresivamente la información mostrada al modelo a medida que avanza el entrenamiento, buscando favorecer la exploración inicial y disminuir luego la dependencia de esas ayudas. Esta retirada afecta a los criterios incluidos en la consulta, mientras que la rúbrica completa se conserva para calcular la recompensa mediante un modelo evaluador que comprueba el cumplimiento de cada criterio. En sus experimentos con Qwen2.5-7B-Instruct, los autores evaluaron sobre HealthBench-500, un subconjunto de 500 casos seleccionados aleatoriamente de HealthBench, y utilizaron los 4.500 restantes para entrenamiento, una partición que permite identificar el conjunto original del benchmark, aunque el artículo no lo denomina expresamente así frente a HealthBench Hard y HealthBench Consensus. Bajo ese protocolo, el puntaje pasó de 23,4 para el modelo inicial a 56,4 con RuscaRL, frente a 52,0 con entrenamiento por refuerzo basado en rúbricas sin el mecanismo de ayudas. Para esta evaluación, los autores utilizaron GPT-4.1 como juez, con el formato de evaluación por criterio de HealthBench, y reportaron el promedio de tres ejecuciones en escala porcentual \cite{zhou2026ruscarl}. Aunque describen la suma ponderada y la normalización de los criterios para la recompensa de entrenamiento, no explicitan por separado la fórmula completa del puntaje global del benchmark ni el tratamiento de los valores negativos, por lo que esa descripción no permite confirmar que reproduzcan íntegramente el cálculo oficial. Esta comparación permite examinar el aporte del andamiaje más allá del uso de rúbricas como recompensa, aunque debe considerarse que el entrenamiento médico también utilizó casos de HealthBench distintos de los reservados para evaluación. Así, el trabajo ofrece un antecedente para retirar progresivamente las ayudas, pero queda abierta su aplicación cuando los datos de entrenamiento provienen de un corpus clínico independiente del benchmark.

